\SourceLeaf{275}
\[
\begin{aligned}
Q+q&=46^0\ 40'\ 37''69794\\
Q+q&=58\quad39\quad44{,}09285
\end{aligned}
\]
die endliche Formel dagegen
\[
\begin{aligned}
Q+q&=46^0\ 40'\ 37''69794\\
Q+q&=58\quad39\quad44{,}09283
\end{aligned}
\]
Die Ergebnisse stimmen also so genau überein, wie es zehnstellige Logarithmen erlauben.

\Article{7}
\indent Auf ähnliche Weise lässt sich der Logarithmus von \(m\) in eine Reihe entwickeln. Ihre ersten Glieder sind
\[
\begin{aligned}
\log m={}&-\frac{\sin2\varphi^2}{6\cos\theta^4}\mathord{.}csp^3-\frac{\sin2\varphi^2}{24\cos\theta^6}(cc+11eess)p^4\\
&+\frac{\sin2\varphi^2}{120\cos\theta^8}\mathord{.}\frac{s}{c}(2cc-3ss-ee(40c^4-20ccss-6s^4)\\
&\qquad-e^4ss(104c^4+22ccss+3s^4))p^5
\end{aligned}
\]
Auch das folgende Glied habe ich (auf einem anderen Weg) entwickelt, allerdings nur bis zum Hauptteil seines Koeffizienten, der von der Ordnung \(ee\) ist. Ich erhielt
\[
+\frac{\sin2\varphi^2}{720\cos\theta^{10}}\mathord{.}\frac{1}{cc}(2c^4-18ccss-15s^4)p^6
\]
\indent Der Logarithmus, den diese Reihe ausdrückt, ist der hyperbolische Logarithmus, und \(p\) wird wie oben in Teilen des Radius angegeben. Will man den Logarithmus nach Briggs, wobei \(p\) Grad bezeichnet, so muss noch der Modul als Faktor hinzukommen und \(\tfrac{\pi p}{180}\) anstelle von \(p\) geschrieben werden. In dieser Form ergibt sich für unser Beispiel
\[
\begin{aligned}
\log m={}&-0{,}0049612433(\frac{p}{100})^3\\
&-0{,}0017329876(\frac{p}{100})^4\\
&-0{,}002393772(\frac{p}{100})^5\\
&-0{,}0124746(\frac{p}{100})^6
\end{aligned}
\]
